\chapter{Strong C2 pitchfork with arbitrary complementary scales}
\label{component:pf7943b34a7dc}
\noindent\textbf{Component }\path{pf7943b34a7dc}\par
{\small\noindent Source: \path{R13/baseline/baseline/baseline/baseline/advances/c2_strong/strong_complement_pitchfork.tex}\par}
{\small\noindent SHA256: \path{f7943b34a7dce78326e3e9ae90456999195c03136f243662c3381e74e9e416bb}\par}
{\small\noindent Inventory use: SC1--SC23 unbounded C or alpha and all upper event scales, full three normals, strict real boundary Schur, complete cap/corner outer proof\par}
\medskip
\section[Strong C2 pitchfork with arbitrary complementary scales]{Strong C2 pitchfork with arbitrary complementary scales}
\label{pf7943b34a7dc:sec:strong-complement-pitchfork}

\begin{theorem}[No upper bound on the two complementary events]
\label{pf7943b34a7dc:thm:strong-complement-pitchfork}
Let $K_c,K_t\supset\mathbb R$ be independent real closed fields.
Let the original canonical C2 parent be physical over $K_c$, with
$\lambda=\beta=0$, $d$ a positive bounded unit, $C,\alpha>0$
having bounded reciprocals, and $n,m,A,D\geq0$ arbitrary.
Suppose it has an actual finite noncommuting stationary zero with
$u_c=k>0$ infinitesimal and $v_c=a>0$ a bounded unit.
Then the entire original parent is tensor SOS at every independent
projective measurement pair and every real or complex projective state.
There is no upper bound or relative-scale condition on any event weight.
\end{theorem}

The case of bounded $C,\alpha$ is the existing three-large-event
theorem. We prove both remaining cases, without modifying the parent
or identifying the source and testing fields. We use
\[
\begin{gathered}
 D_u=1+u^2,\quad D_v=1+v^2,\quad r=uW-Y,\quad p=vY-Z,\\
 J=uW+uvX-Y-vZ,\quad R_2=vW-X+up,\quad B=X+uZ,\\
 H=2\{J^2+D_up^2+D_vB^2-D_uD_vZ^2\},\\
 F=dH+4nD_uD_vX^2+4mR_2^2+4CD_vr^2
       +4AD_uD_vY^2+4DJ^2+4\alpha D_up^2.              
\end{gathered}
\tag{SC1}
\]
Every square in this formula retains its original coefficient.

\paragraph{Necessary reductions from the full parent.}
Cap only $C,\alpha$ at positive bounded units in the initial hard
state estimate. At the actual angles the bounded selected form is
bounded below by
$b(X^2+Y^2+Z^2)-B_0k^2W^2$, with positive real $b,B_0$.
Every retained packet is nonnegative. Hence the contact has $W_c\ne0$;
on setting $W_c=1$,
\[
 X_c,Y_c,Z_c=O(k),\qquad m=O(k^2).
\]
Each positive event packet is $O(k^2)$ as well. Thus, writing
$y=Y_c/k$, $z=Z_c/k$, $x_1=X_c/k$, $j=J_c/k$,
\[
 Ay^2,\ Dj^2,\ C(1-y)^2,\ \alpha(ay-z)^2=O(1),\qquad
 j=1-y-az+aX_c.                                      \tag{SC2}
\]
If $C$ is unbounded then $\bar y=1$, so $A$ is bounded. If $D$
were also unbounded, $j,z$ would be infinitesimal, and the residue of
\[
 H_c/(2k^2)=j^2+(1+k^2)(ay-z)^2
            +(1+a^2)(x_1+Z_c)^2-(1+k^2)(1+a^2)z^2
\]
would be $\bar a^2+(1+\bar a^2)\bar x_1^2>0$.
This contradicts $H_c\leq0$ at the zero of (SC1). Therefore
$C$ unbounded forces both $A,D$ bounded. If $\alpha$ is unbounded
and $A$ were unbounded, (SC2) would make $y,z$ infinitesimal and
$\bar j=1$; the same residue would instead be
$1+(1+\bar a^2)\bar x_1^2>0$. Thus $\alpha$ unbounded forces $A$
bounded. It remains to prove precisely
\[
 \text{(I) } C\text{ unbounded},\ A,D\text{ bounded};\qquad
 \text{(II) } C\text{ bounded},\ \alpha\text{ unbounded},\ A\text{ bounded}.
                                                               \tag{SC3}
\]

\paragraph{Uniform stationary scale entry.}
Set $U=k^2$, $V=a^2$, $N=d+2n$, $\theta=n/N$, and
\[
\begin{gathered}
 t_2=d+2(\alpha+m+A),\quad T=t_3=d+2(C+D),\quad
 t_4=d+2(\alpha+m),\\
 t_6=d+2(\alpha+C+A),\quad t_8=d+2(C+A+D).
\end{gathered}
\]
The exact original Laurent stationary equations are
\[
\begin{split}
 U^2 A t_4\{\theta T+C(1+2D/N)V\}
     &=mV\{\alpha t_8+Dt_6V\},\\
 \alpha A(1+2m/N)T+\theta A Tt_4U
     &=\{mDt_6/U+CDt_2\}V^2.                         
\end{split}\tag{SC4}
\]
Define only pure coefficient quantities
\[
 \delta=CD/T,\quad e=(D/T)(t_6/\alpha),\quad b=t_2/\alpha,
 \quad \mathcal B=\theta+(C/T)(1+2D/N)V,
 \quad \mathcal G=t_8/T+eV.
\]
Because one of $C,D$ is bounded in either case (SC3),
$\delta\leq\min(C,D)/2$ is bounded. The identities
\[
 e=(D/T)\{2+(d+2A)/\alpha\}+2\delta/\alpha,
 \quad \mathcal G=1+2A/T+eV,
 \quad \mathcal B=\theta+(C/T)V+(2\delta/N)V
\]
show boundedness of all these quantities. The quantities
$b,\mathcal G,t_4/\alpha$ are positive bounded units. So is
$\mathcal B$: since $C,N$ have bounded reciprocals and $d$ is
bounded,
\[
 \frac{T}{C(N+2D)}\leq\frac d{CN}+\frac2N+\frac1C=O(1),
\]
and consequently
\[
 \theta+(C/T)(1+2D/N)
 =\frac{n+C(N+2D)/T}{d+2n}
 \geq\frac{n+c_0}{d+2n}
\]
for a positive real $c_0$. The bounded-unit factor $V$ preserves
this lower bound. Dividing (SC4) by $\alpha T$ gives exactly
\[
 \frac m{U^2}=A\rho,\qquad
 A\mathcal H=\delta bV^2,\qquad
 \rho=\frac{(t_4/\alpha)\mathcal B}{V\mathcal G},\quad
 \mathcal H=1+2m/N+\theta(t_4/\alpha)U-U\rho eV^2.      \tag{SC5}
\]
Here $\rho$ is a positive bounded unit and $\mathcal H=1+O(U)$.
These identities were obtained before inverting $A$ or $D$.
They show $A\asymp\delta$. Since $C$ has a bounded reciprocal,
$\delta$ infinitesimal implies $D$ infinitesimal.
If $A$ were infinitesimal, therefore, so would $D,m$ be.
But the exact physical test
\[
 u=v=1,\qquad (W,X,Y,Z)=(1,0,1,1)
\]
has $X=r=p=0$, $R_2=1$, $J=-1$, $H=-2$, and hence
$F=-2d+4m+16A+4D$. Its residue would be negative. We conclude
\[
 A\asymp1,\qquad D^{-1}=O(1),\qquad
 m/k^4\asymp1,\qquad \gamma=4ma^2/k^4\asymp1.        \tag{SC6}
\]
In case (I), $D$ is a bounded unit; in case (II), it may be
arbitrarily large. In both cases $\delta$ is a positive bounded unit.

\paragraph{Full state equations and bounded graph data.}
For later use, the three entire state derivatives of (SC1) are
\[
\begin{split}
 F_X/4={}&D_v[d+2nD_u]X+(d+2D)uvJ+d uD_vZ-2mR_2,\\
 F_Y/4={}&-(d+2D)J+(d+2\alpha)vD_up-2CD_vr
                  +2AD_uD_vY+2muvR_2,\\
 F_Z/4={}&-(d+2D)vJ-(d+2\alpha)D_up+d uD_vB
                  -dD_uD_vZ-2muR_2.                
\end{split}\tag{SC7}
\]
In case (I), the $Z$ equation first yields $\alpha p_c=O(k)$,
because $D$ is bounded and every core coordinate is $O(k)$.
Then the $Y$ equation yields $Cr_c=O(k)$. In case (II), $C,A$
are bounded, so the $Y,Z$ equations are a system with bounded right
side $O(k)$ and unit coefficient matrix
\[
 \begin{pmatrix}-1&a\\-a&-1\end{pmatrix}
\]
on $((d+2D)J_c,(d+2\alpha)(1+k^2)p_c)$.
They yield $DJ_c,\alpha p_c=O(k)$. Finally the $X$ equation in
either case yields $NX_c=O(k^2)$. Thus put $x=X_c/k^2$ and
\[
\begin{array}{ll}
 \text{(I)}&r_0=r_c/k,\ p_0=p_c/k,
       \qquad Nx,Cr_0,\alpha p_0=O(1),\\
 \text{(II)}&j_0=J_c/k,\ p_0=p_c/k,
       \qquad Nx,Dj_0,\alpha p_0=O(1).
\end{array}                                                    \tag{SC8}
\]
All unweighted graph constants are bounded as well. The reciprocals
of $N,C,D,\alpha$ used below are bounded pure coefficient elements.

\paragraph{Two exact graphs retaining the large packets.}
For independent finite testing angles, case (I) uses
\[
 X_g=xu^2W,\quad r_g=r_0uW,\quad p_g=p_0uW,
 \quad Y_g=(1-r_0)uW,\quad Z_g=\{v(1-r_0)-p_0\}uW.    \tag{SC9}
\]
Its exact normals and inverse are
\[
\begin{gathered}
 E_X=\sqrt N(X-xu^2W),\quad E_r=\sqrt C(r-r_0uW),
 \quad E_p=\sqrt\alpha(p-p_0uW),\\
 X=xu^2W+E_X/\sqrt N,\quad
 Y=(1-r_0)uW-E_r/\sqrt C,\quad
 Z=vY-p_0uW-E_p/\sqrt\alpha.                         
\end{gathered}
\tag{SC10}
\]
Case (II) uses
\[
 X_g=xu^2W,\quad J_g=j_0uW,\quad p_g=p_0uW,
 \quad Y_g=\frac{u(1-j_0+vp_0+xu^2v)}{D_v}W,
 \quad Z_g=vY_g-p_0uW.                              \tag{SC11}
\]
Here the exact normals and inverse are
\[
\begin{gathered}
 E_X=\sqrt N(X-xu^2W),\quad E_J=\sqrt D(J-j_0uW),
 \quad E_p=\sqrt\alpha(p-p_0uW),\\
 X=xu^2W+E_X/\sqrt N,\quad J=j_0uW+E_J/\sqrt D,
 \quad p=p_0uW+E_p/\sqrt\alpha,\\
 Y=(uW+uvX+vp-J)/D_v,\qquad Z=vY-p.                 
\end{gathered}
\tag{SC12}
\]
Both graphs match the entire actual contact. All denominators in
these formulas are pure coefficient factors or the pure testing
unit $D_v$. No inverse of a mixed tensor element occurs.
The possibly unbounded packets become respectively
\[
\begin{array}{ll}
 \text{(I)}&4(n/N)D_uD_vE_X^2,quad4D_vE_r^2,quad4D_uE_p^2,\\
 \text{(II)}&4(n/N)D_uD_vE_X^2,quad4E_J^2,quad4D_uE_p^2
\end{array}
\]
in the quadratic normal remainder. Their graph and first-normal
coefficients are bounded by (SC8), for example
$nx,nx^2,Cr_0^2,\sqrt C r_0,\alpha p_0^2,\sqrt\alpha p_0$
in case (I), and $Dj_0^2,\sqrt D j_0$ in case (II).
Every remaining source coefficient is bounded.

\paragraph{Entire scalar and entire first-normal division.}
Set $W=1$, $K=k^2$, and
$R=u^2-(K/a)v$, $T_s=u(v-a)$.
Since $R_2|_g=v+(p_0-x)u^2$ for both graphs, the entire scalar has
the exact shape
\[
 \Phi=F|_g=4mv^2+u^2H_0(u^2,v),                    \tag{SC13}
\]
with bounded polynomial coefficients in case (I), and bounded
coefficients in $B_c[v,D_v^{-1}]$ in case (II).
The three scalar contact jets give
$H_0(K,a)=-\gamma K$, $H_{0,U}(K,a)=\gamma$, and
$H_{0,v}(K,a)=-2\gamma K/a$. In particular the explicit expression
\[
 \Psi=H_0-\gamma u^2+2\gamma(K/a)v
\]
satisfies $\Phi-\gamma R^2=u^2\Psi$, and its value and both first
derivatives vanish at $(u^2,v)=(K,a)$. Finite Taylor division of
$\Psi$ therefore gives
\[
 \Phi=\gamma R^2+u^2\{(u^2-K)^2H_{20}
       +(u^2-K)(v-a)H_{11}+(v-a)^2H_{02}\},           \tag{SC14}
\]
with bounded coquotients. This is an exact bounded Gram on
$(R,T_s)$, with separate residue pivot $\bar\gamma$ on $R$.
For the rational case, Taylor division is legitimate because
\[
 D_v^{-r}-D_a^{-r}
 =-(v-a)(v+a)
   \frac{\sum_{i=0}^{r-1}D_v^iD_a^{r-1-i}}{D_a^rD_v^r}.
\]
Iterating this identity gives the second-order division; only pure
factors have been inverted.

Here are the full first normals, rather than initial forms. In
case (I) they are
\[
 \mathcal N_X=F_X|_g,\quad
 \mathcal N_r=-(F_Y+vF_Z)|_g,\quad
 \mathcal N_p=-F_Z|_g,                              \tag{SC15}
\]
and in case (II) they are
\[
 \mathcal N_X=\{F_X+(uv/D_v)(F_Y+vF_Z)\}|_g,\quad
 \mathcal N_J=-(F_Y+vF_Z)|_g/D_v,\quad
 \mathcal N_p=(vF_Y-F_Z)|_g/D_v.                     \tag{SC16}
\]
The derivatives in (SC7) specify every term. In particular the
following exact packet expressions expose the cancellations:
\[
\begin{split}
 -(F_Y+vF_Z)/(4D_v)
   &=(d+2D)J+2Cr-2AD_uY-duvB+dvD_uZ,\\
 (vF_Y-F_Z)/(4D_v)
   &=(d+2\alpha)D_up-2Cvr+2AvD_uY+2muR_2-duB+dD_uZ.
\end{split}\tag{SC17}
\]
The actual coefficients of the normal roots are obtained by dividing
these rows by $\sqrt N,\sqrt C,\sqrt\alpha$ in case (I), and by
$\sqrt N,\sqrt D,\sqrt\alpha$ in case (II).
Using (SC8), their exact forms after those divisions are
\[
 \mathcal N_X/\sqrt N=-8mv/\sqrt N+u^2L_X(u^2,v),
 \qquad \mathcal N_i/\sqrt{w_i}=uL_i(u^2,v),          \tag{SC18}
\]
where all $L_i$ have bounded coefficients in the indicated ring.
The contact equations give
$L_X(K,a)=8ma/(K\sqrt N)$ and $L_i(K,a)=0$.
Writing $L_X=L_X(K,a)+(u^2-K)L_{X1}+(v-a)L_{X2}$, and similarly
for $L_i$, gives the full exact row identities
\[
\begin{split}
 \mathcal N_X/\sqrt N
  &=(8ma/(K\sqrt N)+u^2L_{X1})R
       +u\{(K/a)L_{X1}+L_{X2}\}T_s,\\
 \mathcal N_i/\sqrt{w_i}
  &=uL_{i1}R+\{(K/a)L_{i1}+L_{i2}\}T_s.             
\end{split}\tag{SC19}
\]
Thus the $X$ row has infinitesimal couplings to both soft roots,
and the other rows have infinitesimal couplings to $R$.
The exact state degree is two, so the bounded quadratic remainder
already accounts for every higher-normal term. Equations
(SC10)--(SC19) give an exact full-parent bounded five-root Gram.

\paragraph{A common exact real Schur certificate.}
The $X$ and $R$ residue pivots are $2(1+\bar a^2)$ and
$\bar\gamma$. The unscaled boundary hard form on $Y,Z$ is
\[
 \mathsf B(v)[Y,Z]
 =sD_vY^2+4D(Y+vZ)^2+4\alpha(vY-Z)^2,
 \qquad s=4(C+A)+2d.                               \tag{SC20}
\]
In case (I), $Y=-E_r/\sqrt C$,
$Z=-vE_r/\sqrt C-E_p/\sqrt\alpha$ at $u=0$.
The form has the explicit pivots $4D_vE_r^2+4E_p^2$ plus a
bounded nonnegative form. In case (II),
$Y=(vE_p/\sqrt\alpha-E_J/\sqrt D)/D_v$ and
$Z=(-E_p/\sqrt\alpha-vE_J/\sqrt D)/D_v$.
The explicit pivots are $4E_J^2+4E_p^2$, again with a bounded
nonnegative remainder. The hard residue is strict in both cases.

The determinant and the exact boundary Schur numerator are
\[
\begin{gathered}
 \det\mathsf B(v)=4(1+v^2)\{s(D\,v^2+\alpha)+4D\alpha(1+v^2)\},\\
 Q_{10}=4\alpha A[d+2(C+D)],\qquad
 Q_{12}=4CD[d+2(\alpha+A)],\\
 Q_0=8ACD+8AC\alpha-4ADd-2Ad^2-2Cd^2-4Cd\alpha
       -2Dd^2-4Dd\alpha-d^3-2d^2\alpha.             
\end{gathered}
\tag{SC21}
\]
Put $S=\alpha T$. Every coefficient of $\det\mathsf B/S$ and of
$Q_{10}/S,Q_0/S,Q_{12}/S$ is bounded. This follows directly from
(SC21), bounded $A$, bounded $C/T,D/T,\delta$, and bounded
$1/\alpha$. Moreover
\[
 Q_{10}/S=4A\asymp1,\qquad
 Q_{12}/S=4\delta[d+2(\alpha+A)]/\alpha\asymp1.       \tag{SC22}
\]
The normalized determinant has a positive lower bound for every
bounded unit real $v$. Indeed it equals $C/T$ times the determinant
of the scaled hard form in case (I), and equals
$(D/T)D_v^2$ times that determinant in case (II). The factors
$C/T$ and $D/T$, respectively, are positive units in their cases.

We now take residues of these bounded coefficients and use an
ordinary real variable $v$. The exact real hard Schur function is
\[
 \bar f(v)=
 \frac{8(1+v^2)}{\overline{\det\mathsf B/S}(v)}
 \left\{\overline{Q_{10}/S}+\overline{Q_0/S}v^2
                      +\overline{Q_{12}/S}v^4\right\}.          \tag{SC23}
\]
This formula is an ordinary real Schur calculation, never an
inverse of $\det\mathsf B(v)$ in $K_c\otimes K_t$.
The contact first-normal identities identify the graph at $\bar a$
with the unique minimizer of the scaled real hard form. The scalar
contact jets give $\bar f(\bar a)=\bar f'(\bar a)=0$.
Because $\bar a\ne0$, the quartic bracket in (SC23) is exactly
\[
 \overline{Q_{12}/S}(v^2-\bar a^2)^2.
\]
Its strict second derivative supplies the soft Schur pivot
\[
 \tfrac12\bar f''(\bar a)=
 \frac{32(1+\bar a^2)\bar a^2\overline{Q_{12}/S}}
      {\overline{\det\mathsf B/S}(\bar a)}>0.
\]
The entire five-root Gram therefore has strictly positive residue
when $u$ is infinitesimal or zero and $\bar v=\bar a$.
The strict-residue Gram lemma proves tensor SOS. The exact identity
\[
 F(u,-v;W,-X,Y,-Z)=F(u,v;W,X,Y,Z)
\]
supplies $\bar v=-\bar a$.

\paragraph{All other independent tests, including literal endpoints.}
If $u$ is infinitesimal and $v$ is a bounded unit with
$\bar v\ne\pm\bar a$, retain the full original $m,n$ packets.
In case (I) use the roots
\[
 X,\quad K_C=\sqrt C r,\quad K_\alpha=\sqrt\alpha p,\quad uW
\]
and the exact inverse
$Y=uW-K_C/\sqrt C$, $Z=vY-K_\alpha/\sqrt\alpha$.
In case (II) use
\[
 X,\quad K_D=\sqrt D J,\quad K_\alpha=\sqrt\alpha p,\quad uW
\]
and the inverse in (SC12), with unshifted $J,p$.
The entire remaining parent is a bounded quadratic Gram on these
four roots. Its separate $X$ residue pivot is
$2\bar d(1+\bar v^2)$; its hard pair is the strict form (SC20)
in the scaled coordinates just displayed, and its soft Schur
complement is $\bar f(\bar v)>0$ by (SC23).

For the two endpoint neighborhoods $u,v\sim0$ and
$u,q=1/v\sim0$, divide by $d$, choose
\[
 \mu=\min(1,C/d,A/d,D/d,\alpha/d)>0,
\]
a bounded unit by (SC6), and retain all excesses together with
the full $m,n$ packets. The selected bounded parent has signed
weight one, $C=A=D=\alpha=\mu$, and $m=n=0$.
In the finite chart it is exactly a Gram on $(X,Y,Z,uW)$.
At $(u,v)=(0,0)$ its hard block is positive and its soft Schur
complement is $8Q_{10,\mu}/\det\mathsf B_\mu(0)>0$.
For the reciprocal chart use the polynomial homogeneous parent
$q^2F_\mu(u,1/q)$, on the same four roots. At $(u,q)=(0,0)$
its hard block is
$\operatorname{diag}(2,4(C+A+\alpha)+2,4D)$ in the
$(X,Y,Z)$ variables, evaluated at the capped weights, and its
soft Schur complement is
$8Q_{12,\mu}/\det\mathsf B_{\mu,\infty}>0$.
These are bounded polynomial Grams, so they include the full
infinitesimal neighborhoods and literal $v=0$ or $q=0$.

For completeness, the other open commuting axes use the same cap;
here the other angle has a nonzero finite residue.
At $v=0$ the core is $2\{(uW-Y)^2+(X+uZ)^2\}$ and the
$A,\alpha$ packets force $Y=Z=0$; finite nonzero $u$ then forces
$W=X=0$. At $u=\infty$ the homogeneous core is
$2\{(W+vX)^2+(vY-Z)^2\}$; the $C,A$ packets force $W=Y=0$,
and finite nonzero $v$ forces $X=Z=0$. At $v=\infty$ the core
is $2(1+u^2)(X^2+Y^2)$; the $C,D$ packets force $W=Z=0$
for finite nonzero $u$. Thus these complete state matrices are strict.

Only the two corners with $u=\infty$ remain. Put $p_u=1/u$.
At $(p_u,v)=(0,0)$ the selected state form is
$(2+8\mu)W^2+4\mu Y^2+(2+4\mu)Z^2$.
Its sole zero ray is $X$. In the chart $X=1$, the five small
variables $(p_u,v,W,Y,Z)$ have exact quadratic part
\[
 (2+4\mu)(W+v)^2+2(Z+p_u)^2+4\mu(W^2+Y^2+Z^2)\succ0.
\]
At $(p_u,q)=(0,0)$ the selected state form is
$4\mu W^2+(2+4\mu)X^2+(2+8\mu)Y^2$.
In its sole zero-ray chart $Z=1$, the variables
$(p_u,q,W,X,Y)$ have exact quadratic part
\[
 2X^2+(2+4\mu)(Y-q)^2+4\mu\{W^2+Y^2+(X-p_u)^2\}\succ0.
\]
There are no lower-order terms. Allocating every higher monomial
to two coordinate roots leaves an infinitesimal factor, giving a
strict-residue Gram for the entire selected polynomial at each
corner. Retain the original excess packets throughout.
After pure testing projective normalization, states outside each
displayed zero ray have positive residue value and are covered by
the scalar strict-residue lemma.

Noncommuting-residue tests use the inherited arbitrary-physical-source
theorem. This exhausts all independent projective tests. Finally
homogenize the local graph identities in $W$, replacing $R,T_s$
by $RW,T_sW$; they then include $W=0$. Restore all pure positive
coefficient and testing factors, and split complex states into real
and imaginary parts. No source infinitesimal was compared with an
independent testing infinitesimal anywhere in the argument.

\paragraph{Exact scope.}
This completes the strong C2 finite pitchfork for $C/d,\alpha/d$
bounded below, with arbitrary upper scales of all six events.
The reflected actual-contact domain follows by the established C2
reflection. It does not cover other actual-angle residue types or
sources in which either complementary ratio has no positive real
lower bound; no exhaustive global C2 join is asserted here.
